%    Latex Template for ACA 2016
%%%%%%%%%%%%%%%%%%%%%%%%%%%%%%%%%%%%%%%%
%    DO NOT CHANGE THE SETTINGS BELOW
%%%%%%%%%%%%%%%%%%%%%%%%%%%%%%%%%%%%%%%%
\documentclass[11pt]{article}
\usepackage{epsfig}
\usepackage{graphicx}
\usepackage[latin1]{inputenc}
\usepackage[T1]{fontenc}
\usepackage{mathptmx}
\usepackage{url}



\begin{document}



%%%%%%%%%%%%%%%%%%%%%%%%
%%%%%%%%%%%%%%%%%%%%%%%%%%%%%%%%%%%%%%%%
%    DO NOT CHANGE THE SETTINGS ABOVE
%%%%%%%%%%%%%%%%%%%%%%%%%%%%%%%%%%%%%%%%
%
\vspace{-0.8cm}
\begin{flushleft}
\Large \textbf{\noindent
%%% PLEASE INSERT THE TITLE OF YOUR TALK HERE %%%%%%%
Developing Competences in Higher Mathematics
in a CAS Supported Learning Environment
}
%%%%%%%%%%%%%%%%%%%%%%%%%%%%%%%%%%%%%%%%%%%%%%%%%%%%%%%%
\\
%%%%%%%%%%%%%%%%%%%%%%%%%%%%%%
\vspace{0.5cm}
\normalsize
 %%%%%%             NAMES OF AUTHORS         %%%%%%
 %%%%%% NAME OF SPEAKER SHOULD BE UNDERLINED %%%%%%%%%
\normalsize{
 %%%%%% FOR EXAMPLE %%%%%
 \underline{E. Varbanova}$^1$, M. Durcheva$^2$
 %%%%%%%%%%%%%%%%%%%%%%%%
} \\
\vspace{5mm}
\textit{\footnotesize
 %%%%%% AFFILIATION OF AUTHORS %%%%%%
$^1$ Technical University of Sofia, Bulgaria,
elvar@tu-sofia.bg\\
$^2$ Technical University of Sofia, Bulgaria,
mdurcheva66@gmail.com\\
 %%%%%%%%%%%%%%%%%%%%%%%%%%%%%%%%%%%%%
}
\end{flushleft}

\vspace{5mm}

\textbf{1.	Competency-based Objectives as a Challenge to the University Culture}

\vspace{5mm}

At the Technical University of Sofia the teaching and learning of Higher mathematics aims at the synergy of two modes of learning: traditional mode and technology supported mode. The effect of synergy of their advantages as well as the danger of synergy of their drawbacks need attention. Any combination of these two modes is referred to as blended learning. There exist many different versions of their combination as a result of teachers' creativity in the design and development of learning scenarios. Harmonized and successful scenarios will be further produced by teachers: it takes time and requires experience.


The term 'competence or competency' is widely used in varying ways. The use depends on the context: competence has different emphasis in business processes than in formal education organizations or in theoretical discussion. The term 'key competences' appears in the guidelines of the Bologna Process, the European frame of reference for the reform of higher education (~\cite{CHUR1, CHUR2}). In Refs.~\cite{TRACE} the principal practices and approaches with regard to the development and assessment of individual competencies have been reported.  In the light of the current and future impact of this rapidly evolving field on the Competency Frameworks , it is expected mathematics teachers to update the principles and approaches in mathematics education in order to reflect the latest evolutions. Various frameworks have been developed to aid in the development of competencies outlining what graduates should know and be able "to do" as a result of their education ~\cite{FORD}.

    The term 'competency' aims at the identity of the student and his/her deliberate action while dealing with complex tasks and challenges. It has multiple layers and is applied as competence-in-something, including both cognitive and non-cognitive dimensions ~\cite{CHUR1}. Its cognitive meaning designates subject-related and general cognitive abilities up to comprehension and problem-solving, whereas its action-based meaning contains cognitive but also motivational, emotional, social and value-related components. Key competences ~\cite{CHUR2} encompass generative, context-independent abilities and are closely related to action competences.

  \vspace{5mm}


\textbf{	2.	Developing Competences in Higher Education}

\vspace{2mm}

Competences are part of today's curricula in higher education. They are the objective of a contemporary education: some are explicit learning objectives, others are implicitly contained in studying. The re-orientation of the \emph{curricula towards outcome and competency} makes students autonomous responsible subjects in the teaching-learning-assessment (TLA) process. TLA have to go beyond declarative knowledge and skills and lead to \emph{procedural competences of problem solving}. Both teachers and students face each other in new roles in the TLA process: \emph{the teachers as 'facilitators'} - demanding and encouraging the students, \emph{the students as active partners}. Both sides purposefully work to achieve the educational objectives.

The shift from teaching to learning provides for a frame of orientation and development, that it brings to the fore the freedom of the individual learner that is inherent in the studying process. Competence-based objectives assure more freedom compared to knowledge-based goals.  As an ability of handling knowledge, competences can be seen as an open integral of various possible learning contents.

  \vspace{2mm}


\textbf{3. A Combined Learning Environment for the Acquisition of Competences in Higher Mathematics}

\vspace{2mm}

The integrative and additive support with CAS incorporated in blended learning can help lead the students to a connection with subject-specific competences. Lectures, seminars and laboratory classes in Mathematics are to connect key competences with mathematics-related competences through extended learning environments.
Depending on the assessment model CAS can often come to the rescue - the student can use the entire potential of CAS during the semester for different purposes of each assessment tasks in both models: continuous assessment and formative assessment with elements of continuous.

In order to be effective, a combined learning environment is to aim at competences beyond knowledge: self-exploring and self-reflecting as well as being opened up for communication and feedback.

As a demonstration of the above mentioned issues the students' work on a set of tasks - guided by teachers, will be presented.  One such task suitable for individual task, independent learning or coursework is the following:

\textbf{Task 1.} Determine the values of the real parameter $m\in \mathbf{R} $  for which the equation

\begin{equation}
f(x) = (2-m)x^3-3mx^2-3mx+ 2-m = 0 \label{eq1}
\end{equation}

has a double root.
This type of tasks could bring to the fore the student conscious activity. Learning and assessment activities as well as competences and key competences (such like Life Long Learning) will be discussed.

  \vspace{2mm}


\textbf{4. Conclusion }

\vspace{2mm}

There is "no way back", i.e. no education, progress and development without technology. Digital mathematics exists and digital resistance is not appropriate.

%%%%%%%%%%%%%%%%%%%%%%%%%%%%%%%%%%%%%%%


%%%%%%%%%%%%%%%%%%%%%%%%%%%%%%%%%%%%%%%%
%% BIBLIOGRAPHY
%%%%%%%%%%%%%%%%%%%%%%%%%%%%%%%%%%%%%%%%
\begin{thebibliography}{00}
\addcontentsline{toc}{chapter}{References}
\setlength{\itemsep}{-1mm}
\footnotesize


\bibitem{CHUR1} D. Chur, \textit{Developing Key Competences In Higher Education}. In Bauder-Begerow, I. Sch$\ddot{a}$fer, S. (Hrsg.) Learning 9$/$11 - Teaching for Key Competences in Literary and Cultural Studies., Heidelberg: Winter Verlag,  53-74 [PDF-Datei] (2011).
 \bibitem{CHUR2} D. Chur, \textit{Kompetenzorientierung im Studium und der Erwerb von Schlusselkompetenzen}. In: Kossek B. Zwiauer C. (Hrsg.): Universit$\ddot{a}$t in Zeiten von Bologna. Zur Theorie und Praxis von Lehr- und Lernkulturen. Wien: VR unipress, 289-314 [PDF-Datei] (2012).
 \bibitem{FORD}     K. Ford, \textit{Competency-Based Education; History, Opportunities, and Challenges, UMUC Information and Research Analysis Center for Innovation} In Learning and Student Success (CILSS) (2014)   https://www.umuc.edu/innovatelearning/upload/cbe-lit-review-ford.pdf

\bibitem{TRACE} \textit{TRACE Project (TRAnsparent Competences in Europe)}, Overviiew of European Competency Frameworks, http://www.menon.org/wp-content/uploads/2012/11/9.-TRACE-Overview-of-EU-competency-frameworks1.pdf, (2012)



\end{thebibliography}
\normalsize
\end{document}
